\lecture{3.4}{d-Separation in SWIGs}
\label{lec:3-4}

\subsection{Paths in SWIGs}

SWIGs contain both random nodes and fixed intervention nodes. d-Separation in a SWIG is read using paths among the random nodes, excluding the fixed intervention labels such as $a$.

\begin{definition}[Path in a SWIG]
A \emph{path} between two random nodes in a SWIG is a sequence of adjacent random nodes connecting them. Fixed intervention nodes are not included as nodes on the path.
\end{definition}

For identification of the ATE, the key graphical question is often whether
\[
\Ya{a}\ci A\given W
\]
is implied by the SWIG. This is a statement about independence between the random treatment node $A$ and the random counterfactual outcome node $\Ya{a}$.

\subsection{Fixed Intervention Nodes}

The fixed node $a$ represents an intervention setting $A=a$. It is not random and therefore is not conditioned on in the same way as a random variable. It inherits arrows that leave the original treatment node.

\begin{figure}[H]
\centering
\begin{tikzpicture}[node distance=1.8cm,>=Latex]
  \node[draw, circle, fill=blue!5] (W) {$W$};
  \node[draw, circle, fill=green!5, right=of W] (A) {$A$};
  \node[draw, rectangle, fill=gray!12, below=.6cm of A] (a) {$a$};
  \node[draw, circle, fill=orange!8, right=of A] (Y) {$\Ya{a}$};
  \draw[->, thick] (W) -- (A);
  \draw[->, thick] (W) to[bend right=25] (Y);
  \draw[->, thick] (a) -- (Y);
\end{tikzpicture}
\caption{The fixed node $a$ is not part of random-variable paths, but it defines the intervention world.}
\end{figure}

The path from $A$ to $\Ya{a}$ through $W$ is
\[
A\leftarrow W\to \Ya{a}.
\]
The arrow from $a$ to $\Ya{a}$ does not create a random path from $A$ to $\Ya{a}$ because $a$ is fixed.

\subsection{Colliders and Non-Colliders in SWIGs}

As in DAGs, whether a node is a collider depends on the particular path. A random node $V$ is a collider on a path if the two adjacent arrows on that path both point into $V$:
\[
\cdots \to V \leftarrow \cdots.
\]
Otherwise, $V$ is a non-collider on that path.

A node can be a collider on one path and a non-collider on another. Consider paths between $A$ and $\Ya{a}$ in a SWIG with nodes $W,H,Z$. A node $Z$ may be a collider on
\[
\Ya{a}\leftarrow H\to Z\leftarrow W\to A,
\]
but a non-collider on
\[
\Ya{a}\leftarrow Z\leftarrow A.
\]
The path determines the role.

\subsection{Unconditional d-Separation}

A path d-connects two random nodes unconditionally if it has no closed collider and no conditioned non-collider. With no conditioning set, forks and chains are open, while colliders are closed.

Common open unconditional paths include:
\[
X\to \cdots \to Y,
\qquad
X\leftarrow \cdots \leftarrow Y,
\qquad
X\leftarrow \cdots \leftarrow T\to \cdots \to Y.
\]
The third form is a shared common cause path.

In the SWIG
\[
A\leftarrow W\to \Ya{a},
\]
$A$ and $\Ya{a}$ are not unconditionally d-separated because the fork through $W$ is open. This represents confounding by $W$.

\subsection{Conditional d-Separation}

\begin{definition}[Conditional d-separation in a SWIG]
A path between random nodes is open given a conditioning set $S$ if every non-collider on the path is not in $S$, and every collider on the path is either in $S$ or has a descendant in $S$. Two nodes are d-separated by $S$ if every path between them is blocked by $S$.
\end{definition}

This is the same logic as for DAGs, applied to random nodes in the SWIG. Fixed nodes define interventions but are excluded from conditioning sets.

\subsection{Reading $\Ya{a}\ci A\given W$}

The core use of SWIG d-separation in this chapter is to assess conditional randomization. In the confounding SWIG:

\begin{figure}[H]
\centering
\begin{tikzpicture}[node distance=1.8cm,>=Latex]
  \node[draw, circle, fill=blue!5] (W) {$W$};
  \node[draw, circle, fill=green!5, right=of W] (A) {$A$};
  \node[draw, rectangle, fill=gray!12, below=.6cm of A] (a) {$a$};
  \node[draw, circle, fill=orange!8, right=of A] (Y) {$\Ya{a}$};
  \draw[->, thick] (W) -- (A);
  \draw[->, thick] (W) to[bend right=25] (Y);
  \draw[->, thick] (a) -- (Y);
\end{tikzpicture}
\caption{Conditioning on $W$ blocks the backdoor path from $A$ to $\Ya{a}$.}
\end{figure}

The only random path between $A$ and $\Ya{a}$ is
\[
A\leftarrow W\to \Ya{a}.
\]
This is a fork. Without conditioning, the path is open. Conditioning on $W$ blocks the path. Therefore the SWIG implies
\[
\Ya{a}\ci A\given W.
\]

\begin{keyidea}
In a SWIG, conditional randomization is a d-separation statement between the random treatment node $A$ and the counterfactual outcome node $\Ya{a}$.
\end{keyidea}

If an unmeasured variable $U$ also points to both $A$ and $\Ya{a}$, then the path
\[
A\leftarrow U\to \Ya{a}
\]
remains open after conditioning on measured $W$. In that case the graph does not support $\Ya{a}\ci A\given W$.
